\chapter{Glossario}
\label{cap:E-glossario}

\begin{description}[style=nextline, leftmargin=0pt, itemsep=2pt]

\item[\textbf{Ammissione}] La decisione se accettare una richiesta in coda. In
\nf{} \`e presa dal modulo \code{sync-queue} per profondit\`a e per attesa
stimata (\cref{cap:10-sync-queue}).

\item[\textbf{Attesa} (\emph{queue wait})] Intervallo fra l'ammissione della
richiesta e la sua consegna al runtime. Metrica
\code{function\_queue\_wait\_ms}.

\item[\textbf{Budget di concorrenza}] Numero massimo di invocazioni
contemporaneamente in volo su tutte le funzioni in modo \code{BUDGETED}. \`E una
dichiarazione di capacit\`a della piattaforma, non una grandezza inferita dal
carico (\cref{cap:12-concurrency-control}).

\item[\textbf{Ciclo aperto / ciclo chiuso}] Nel primo gli arrivi sono
programmati dall'orologio; nel secondo ogni utente virtuale attende la risposta
prima di inviare. Un carico a ciclo chiuso non pu\`o distinguere politiche di
ammissione (\cref{cap:22-metodo-sperimentale}).

\item[\textbf{Cold start}] Ritardo dovuto all'inizializzazione di un'istanza non
ancora pronta. Metrica \code{function\_cold\_start\_ms}.

\item[\textbf{Compilazione anticipata} (\emph{AOT})] Fase di build di Spring
che risolve i bean e le loro condizioni prima dell'esecuzione. Nell'immagine
nativa le condizioni sulle propriet\`a sono quindi congelate al valore che
avevano in build (\cref{cap:20-build-packaging}).

\item[\textbf{Contratto di saturazione}] Politica comune con cui ogni runtime di
funzione rifiuta il lavoro oltre limiti finiti di conteggio e di byte, con
codici \code{429}, \code{413}, \code{500} e \code{503} e un corpus JSON
condiviso che tutti gli SDK verificano (\cref{sec:18-saturazione}).

\item[\textbf{Control plane}] Il processo che registra le funzioni, orchestra le
invocazioni e governa le politiche. In \nf{} \`e uno solo.

\item[\textbf{Degradazione osservabile}] Riduzione di funzionalit\`a che il
sistema dichiara --- via log, metrica o codice di stato --- anzich\'e
nascondere. \`E il criterio che distingue le degradazioni accettabili da quelle
che fanno rifiutare l'avvio (\cref{cap:03-requisiti-principi}).

\item[\textbf{Deriva della baseline}] Meccanismo per cui un estremo storico
(minimo di latenza, massimo di throughput) viene lentamente riportato verso il
presente, cos\`i che un intervallo anomalo non lo fissi per sempre
(\cref{cap:12-concurrency-control}).

\item[\textbf{Dispatch}] La consegna dell'invocazione al runtime della funzione.

\item[\textbf{Equit\`a max-min pesata}] Criterio di ripartizione in cui si
massimizza l'allocazione minima, chi chiede meno della propria quota la ottiene
per intero, e l'eccedenza \`e ridistribuita~\cite{bertsekas1992datanetworks}.

\item[\textbf{Finestra di ammissione}] La quantit\`a $C + Q$ --- limite di
concorrenza pi\`u dimensione della coda --- oltre la quale una funzione rifiuta.
Il fatto che il limite ne faccia parte \`e il secondo lavoro del limite
(\cref{cap:12-concurrency-control}).

\item[\textbf{Ginocchio}] Il valore del limite di concorrenza a cui il servizio
satura e oltre il quale l'aggiunta produce attesa anzich\'e lavoro. Calcolato
come $\lambda_{\max} S_{\min}$ per la legge di Little.

\item[\textbf{Idempotenza}] Propriet\`a per cui due invocazioni con la stessa
chiave corrispondono a una sola esecuzione (\cref{cap:06-nucleo}).

\item[\textbf{Legge di Little}] $L = \lambda W$: il numero medio nel sistema
uguaglia il tasso di arrivo per il tempo medio di permanenza~\cite{little1961proof}.

\item[\textbf{Limite configurato / effettivo}] Il primo \`e ci\`o che l'utente ha
dichiarato ed \`e un tetto invalicabile; il secondo \`e ci\`o che un controllore
adattivo applica in questo momento (\cref{cap:09-async-queue}).

\item[\textbf{Modulo}] Unit\`a opzionale del control plane, selezionata a tempo
di build e caricata via SPI (\cref{cap:08-sistema-moduli}).

\item[\textbf{Offload}] Inoltro trasparente di un'invocazione sincrona a
un'istanza remota, per politica della funzione o per pressione locale
(\cref{cap:13-offload}).

\item[\textbf{Provider di deployment}] Backend che realizza l'intento
\code{DEPLOYMENT}: Kubernetes, runtime container locale o containerd. I tre
moduli sono mutuamente esclusivi (\cref{cap:15-deployment-provider}).

\item[\textbf{Ricetta}] File YAML versionato (schema v2) che dice quale control plane
costruire --- moduli e modo JVM o nativo --- e quali implementazioni di funzione
impacchettare con esso; non avvia n\'e registra nulla. Quattro comandi Gradle
elencano il catalogo, validano, assemblano e pubblicano
(\cref{sec:20-ricette}).

\item[\textbf{Replay}] Riproduzione immediata dell'esito di un'esecuzione gi\`a
terminata, per una richiesta con chiave di idempotenza corrispondente.

\item[\textbf{Servizio} (\emph{service time})] Intervallo fra dispatch e
completamento, comprensivo del round trip HTTP. Metrica
\code{function\_latency\_ms}. \emph{Non} include l'attesa in coda.

\item[\textbf{Sojourn}] Attesa pi\`u servizio: ci\`o che il chiamante osserva al
netto della rete. Metrica \code{function\_e2e\_latency\_ms}. \`E la grandezza da
cui il modo omonimo decide.

\item[\textbf{SPI} (\emph{Service Provider Interface})] Meccanismo Java di
scoperta di implementazioni via \code{ServiceLoader}, usato per caricare i
moduli.

\item[\textbf{Watchdog}] Binario Rust che fa da entrypoint di ogni container di
funzione gestito e adatta i tre protocolli di runtime a un unico contratto HTTP
(\cref{cap:17-runtime-funzione}).

\end{description}
